El 1921 Albert Einstein (1897-1955) va guanyar el Premi Nobel de Física pel seu treball sobre l'explicació de l'efecte fotoelèctric. Amb el mètode del potencial invers es pot trobar l'energia cinètica màxima dels electrons que són emesos quan s'il·lumina una cel·la fotoelèctrica amb llum monocromàtica. Aquest mètode consisteix a aplicar un potencial entre l'ànode i el càtode per a frenar aquests electrons i impedir que arribin a l'ànode de la cèl·lula fotoelèctrica.

\figura[0.35]{fig1.png}

\begin{apartat}{1,25}
D'acord amb l'esquema de la dreta, a quina placa s'aplica el potencial alt, a l'ànode o al càtode? Quina relació hi ha entre el potencial mínim que cal aplicar per a aturar el corrent (potencial de frenada) i la freqüència dels fotons? Justifiqueu les respostes. En un experiment amb el mètode del potencial invers es van obtenir els resultats indicats a la taula de sota. Determineu la constant de Planck a partir d'aquests resultats.
\end{apartat}

\medskip
{\centering
\begin{tabular}{|l|c|c|c|}
\hline
\textit{Freqüència de la llum incident} ($\times 10^{14}$ \textit{Hz}) & 5,49 & 7,41 & 12,5 \\
\hline
\textit{Potencial de frenada (V)} & 0,40 & 1,19 & 3,29 \\
\hline
\end{tabular}\par}
\medskip

\begin{apartat}{1,25}
Quina és la freqüència llindar? Hi haurà emissió d'electrons si il·luminem la cel·la amb una llum monocromàtica d'una longitud d'ona de 350~nm? Justifiqueu la resposta.
\end{apartat}

\dades{%
  $|e| = 1,602\times 10^{-19}\ \mathrm{C}$.\\
  $c = 3,00\times 10^{8}\ \mathrm{m\,s^{-1}}$.}
